\documentclass[11pt,a4paper]{article}
\usepackage[T1]{fontenc}
\usepackage[utf8]{inputenc}
\usepackage[polish]{babel}
\usepackage{amsmath,amssymb,amsthm,geometry,hyperref}
\geometry{margin=25mm}
\newtheorem{theorem}{Twierdzenie}
\newtheorem{corollary}{Wniosek}
\newcommand{\Tr}{\operatorname{Tr}}
\newcommand{\Dtr}{D_{\rm tr}}
\newcommand{\CQ}{\mathcal C\!\mathcal Q}
\title{FIN: separowalna równowaga i konieczny dysonans kwantowy}
\author{}
\date{}
\begin{document}
\maketitle

\begin{abstract}
Dla kanonicznego kwantowego oddziaływania związanego ze źródłem FIN
konstruujemy jawną separowalną równowagę o obu stanach zredukowanych
równych $C=I/12+W/20$, gdzie $W$ jest zadanym strict.
Stan jest mieszaniną 22 produktów i może być warunkowo przygotowany
przez lokalne pomiary ze wspólnym wyborem podziału, z prawdopodobieństwem
powodzenia $1/2$. Splątanie nie jest więc konieczne w tym przykładzie.
Jednocześnie dowodzimy, że żadna równowaga tego kanonicznego oddziaływania
z marginalnym $C$ nie ma zerowego jednostronnego dysonansu kwantowego.
Wymagane korelacje mogą być separowalne, lecz nie w pełni klasyczne.
Przygotowanie nadal korzysta ze stanu zawierającego dane jądra;
również minimalny plan podziałów nie wyznacza jednoznacznie stanu wspólnego.
Wyniki nie stanowią wyprowadzenia fundamentalnej fizyki FIN.
\end{abstract}

\section{Model, pytanie i rozdzielenie pojęć}

Niech $n=12$, $S$ zamienia dwa nośniki, $D=\sum_i|ii\rangle\langle ii|$
i $|\Omega\rangle=\sum_i|ii\rangle$. Badane oddziaływanie to
\begin{equation}\label{eq:interaction}
 V=\frac{|\Omega\rangle\langle\Omega|+S-2D}{2}.
\end{equation}
Odpowiada ono zadanej mapie źródłowej
$T(\rho)=\Pi\Re\rho$ poprzez $\Tr_2[V(I\otimes\rho)]=T(\rho)$.
Jest to konkretny model mikroskopowy, nie cała klasa możliwych
fizycznych rozszerzeń FIN.
\footnote{Źródła repozytoryjne:
\texttt{FIN\_Quantum\_Source\_Correlation\_and\_Programming\_Report.tex}
oraz \texttt{FIN\_Hartree\_Equivalence\_and\_Strict\_Correlation\_Floor.tex}.
Ostatni raport daje dolną granicę korelacji, lecz nie utożsamia jej
z koniecznością splątania.}

Docelowy stan zredukowany to
\begin{equation}\label{eq:C}
 C=\frac I{12}+\frac W{20},\qquad
 W_{ij}=\frac{\cos((743d_{ij}+650)/4000)}{1+d_{ij}^{9/5}},\quad W_{ii}=0,
\end{equation}
z odległością cykliczną $d_{ij}$. Dane strict pozostają wejściem;
nie wyprowadzamy ponownie ich parametrów.

Stan separowalny ma postać $R=\sum_r p_r\sigma_r\otimes\tau_r$.
Stan klasyczno-kwantowy po stronie pierwszego nośnika ma mocniejszą
postać $Q=\sum_i p_i|a_i\rangle\langle a_i|\otimes\rho_i$, gdzie
$\{|a_i\rangle\}$ jest jedną bazą ortonormalną. Ta druga klasa,
oznaczana przez $\CQ_A$, odpowiada zerowemu jednostronnemu
dysonansowi kwantowemu. Separowalność nie pociąga tej własności.
\footnote{H.~Ollivier, W.~H.~Zurek,
\emph{Quantum Discord: A Measure of the Quantumness of Correlations},
Physical Review Letters 88, 017901 (2001),
\href{https://doi.org/10.1103/PhysRevLett.88.017901}{doi:10.1103/PhysRevLett.88.017901}.
Rozróżnienie separowalności i dysonansu jest znaną treścią literatury,
nie nową definicją FIN.}

\section{Jawna separowalna równowaga strict}

Połóżmy $p=5/11$ i
\begin{equation}\label{eq:program}
 C'=I/12+(C-I/12)/p=I/12+11W/100.
\end{equation}
Dokładne przedziały spektralne dają
\[
 \lambda_{\min}(C')\ge
 \frac{249813284145307}{30000000000000000}>\frac1{125}.
\]
Zatem $C'$ jest dodatnim stanem o jednolitej przekątnej.
\footnote{Przedziały są przeliczane przez
\texttt{fin\_projected\_learning/research.py} i
\texttt{fin\_replication\_consistency/certify.py}.
Nie zastępuje ich dodatniość zaokrąglonej macierzy numerycznej.}

Macierz Hadamarda $H$ rzędu 12 można skonstruować z reszt kwadratowych
modulo 11. Pierwszy wiersz i kolumna są dodatnie, a dolny blok ma
przekątną $-1$ i pozostałe elementy $\chi(j-i)$, gdzie $\chi$ jest
symbolem reszty kwadratowej. Dokładnie $HH^T=12I$.
Jedenaście niepierwszych kolumn daje zrównoważone cięcia po sześć etykiet.
Uwzględnienie obu orientacji daje 22 bloki dodatnie $A_r$ i ich dopełnienia
$B_r$. Każda etykieta należy do 11 bloków $A_r$, a każda para różnych
etykiet do pięciu.

\begin{theorem}[Separowalna równowaga]\label{thm:separable}
Dla
\begin{equation}\label{eq:state}
 \sigma_r=2P_{A_r}C'P_{A_r},\quad
 \tau_r=2P_{B_r}C'P_{B_r},\qquad
 R_{\rm sep}=\frac1{22}\sum_{r=1}^{22}\sigma_r\otimes\tau_r
\end{equation}
stan $R_{\rm sep}$ jest separowalny, ma oba marginały równe $C$
i spełnia $[V,R_{\rm sep}]=0$.
\end{theorem}

\begin{proof}
Każdy czynnik jest dodatni i ma ślad jeden, ponieważ blok zawiera
sześć elementów przekątnej $1/12$. Mamy więc jawny rozkład separowalny.
W marginale przekątna pozostaje $1/12$, a każdy element pozadiagonalny
$C'_{ij}$ jest pomnożony przez $2(5/22)=5/11$. To daje $C$.

Nośniki obu czynników są rozłączne w bazie etykiet, zatem
$DR_{\rm sep}=R_{\rm sep}D=0$ i również
$|\Omega\rangle\langle\Omega|R_{\rm sep}=0$.
Przeciwne orientacje zapewniają $SR_{\rm sep}S=R_{\rm sep}$.
Stąd $[V,R_{\rm sep}]=[S/2,R_{\rm sep}]=0$.
\end{proof}

Nie jest to wnioskowanie ``PPT, więc separowalny''.
Dla rzeczywistego programu rzeczywiście $R_{\rm sep}^{T_B}=R_{\rm sep}$,
lecz separowalność wynika z \eqref{eq:state}.
Ponadto
\[
 R_{\rm sep}\succeq\frac{12}{11}\lambda_{\min}(C')^2(I-D).
\]
Wynika to z sześciu przeciwległych orientacji przypadających na
każdą uporządkowaną parę różnych etykiet. Rząd stanu jest więc
dokładnie 132. Jego prawdopodobieństwo koincydencji
$\Tr(DR_{\rm sep})$ jest zerowe, a dla produktu $C\otimes C$ wynosi $1/12$.

Konstrukcja działa dla dowolnego stanu o jednolitej przekątnej,
jeżeli wzmocniony program jest dodatni. Dla danego rzeczywistego
jądra gwarantowany zakres ma postać
\[
 \gamma\le\frac{5}{11}\frac1{12(-\lambda_{\min}(W))}.
\]
Dla strict granica tego wystarczającego zakresu jest większa niż $0.05$.
Nie jest to dowód maksymalnego zakresu separowalności wszystkich równowag.

\section{Warunkowe przygotowanie przez operacje lokalne}

Wybierzmy równomiernie jedno z jedenastu niezorientowanych cięć.
Na dwóch niezależnych kopiach $C'$ wykonujemy lokalne pomiary
$\{P_A,P_B\}$. Akceptujemy przeciwne wyniki. Krausy sukcesu to
$P_A\otimes P_B$ oraz $P_B\otimes P_A$; dwa pozostałe wyniki
uzupełniają instrument do mapy zachowującej ślad.

Ponieważ $\Tr(P_A C')=\Tr(P_B C')=1/2$, prawdopodobieństwo
powodzenia jest dokładnie $1/2$ dla każdego cięcia. Po zapomnieniu
orientacji i uśrednieniu cięć stanem zaakceptowanej pary jest
$R_{\rm sep}$. Jest to jawna procedura lokalna z klasyczną komunikacją,
bez wejściowego zasobu splątania.

Program $C'$, wspólne losowanie cięcia, porównanie wyników i dostarczanie
kopii są jednak zadanymi zasobami. Program zawiera już macierz $W$.
Wybór zaakceptowanej gałęzi jest jawny i nie stanowi deterministycznej
nieliniowej operacji na jednym nieznanym stanie. Stacjonarność dotyczy
pary po pominięciu rejestru przygotowania; zachowanie rejestru pozwala
rozróżnić niestacjonarne gałęzie warunkowe. Nie udowodniono równowagi
całego układu obejmującego wszystkie rekordy i urządzenia przygotowania.

\section{Dlaczego żadna taka kanoniczna równowaga nie jest klasyczno-kwantowa}

\begin{theorem}[Kryterium odwracalnego bloku]\label{thm:block}
Niech pierwszy marginalny stan stacjonarnego $R$ dla hamiltonianu $H$
ma proste wektory własne $u,v$ z różnymi wartościami $c_u,c_v$.
Jeżeli operator $A=\langle u|H|v\rangle$ na drugim nośniku jest
odwracalny, to $[R,C\otimes I]\ne0$, a zatem $R\notin\CQ_A$.
\end{theorem}

\begin{proof}
Gdyby $[R,C\otimes I]=0$, proste sektory $u$ i $v$ byłyby odłączone
od wszystkich innych sektorów pierwszego nośnika. Oznaczając
$R_u=\langle u|R|u\rangle$ i $R_v=\langle v|R|v\rangle$, blok
$u,v$ równania $[H,R]=0$ daje $AR_v=R_uA$.
Odwracalność wymusza $\Tr R_v=\Tr R_u$, czyli $c_v=c_u$,
sprzeczność. Każdy stan klasyczno-kwantowy komutuje ze swoim
pierwszym marginalnym stanem pomnożonym tensorowo przez $I$.
\end{proof}

Ten ostatni warunek jest jedynie konieczny dla zerowego dysonansu.
Jego spełnienie nie jest certyfikatem klasyczności: na przykład
stan maksymalnie splątany ma marginalny stan proporcjonalny do $I$.
\footnote{C.~A.~Rodríguez-Rosario, G.~Kimura, H.~Imai, A.~Aspuru-Guzik,
\emph{Lazy states: sufficient and necessary condition for zero quantum
entropy rates under any coupling to the environment},
\href{https://arxiv.org/abs/1004.5405}{arXiv:1004.5405}.
Klasa stanów komutujących z marginalnym operatorem jest szersza od
stanów o zerowym dysonansie. Stacjonarność dla jednego $V$ nie oznacza
zerowej odpowiedzi entropii dla każdego innego oddziaływania.}

Dla strict bierzemy rzeczywisty jednorodny wektor $u$ i mod
naprzemienny $v_i=(-1)^i/\sqrt{12}$. Są prostymi wektorami własnymi
kanonicznego $C$ w \eqref{eq:C}, z różnymi wartościami.
Dla $J=\operatorname{diag}((-1)^i)$ mamy dokładnie
\begin{align}\label{eq:cross}
 A&=\frac{|u\rangle\langle v|+|v\rangle\langle u|}{2}-\frac J{12},\\
 A^2&=\frac1{144}(I-P_u-P_v)+\frac{25}{144}(P_u+P_v).
\end{align}
Stąd $\|A^{-1}\|=12$. Twierdzenie \ref{thm:block} dotyczy zatem
\emph{każdej} równowagi kanonicznego $V$ z pierwszym marginalnym $C$,
nie tylko zbudowanego przykładu. Przy obu marginałach równych $C$
dysonans jest niezerowy po obu stronach.

W połączeniu z twierdzeniem \ref{thm:separable} otrzymujemy właściwe
rozstrzygnięcie: splątanie nie jest konieczne, lecz klasyczno-kwantowy
opis korelacji jest wykluczony w tym modelu. Klasyczna komunikacja
w przygotowaniu nie przeczy temu wynikowi, ponieważ lokalne stany
programu zawierają koherencję i nie tworzą jednej wspólnej bazy klasycznej.

\section{Granice ilościowe i ich definicje}

Dla dowolnego $Q\in\CQ_A$, o dowolnym pierwszym marginale $Q_A$,
\[
 \|[R,C\otimes I]\|_1
 \le2(1+\|C\|)\|R-Q\|_1.
\]
Dowód polega na dodaniu i odjęciu $Q$ oraz $Q_A$, użyciu
$[Q,Q_A\otimes I]=0$ i kontrakcji śladu częściowego. Zatem
\begin{equation}\label{eq:distance}
 \inf_{Q\in\CQ_A}\Dtr(R,Q)
 \ge\frac{\|[R,C\otimes I]\|_1}{4(1+\|C\|)}.
\end{equation}

Dla skonstruowanego $R_{\rm sep}$ wszystkie elementy są nieujemne,
a nierówne etykiety mają elementy diagonalne $1/132$.
W elemencie komutatora między $|00\rangle$ i $|10\rangle$ nie ma
kasowania znaków, więc jego moduł jest co najmniej $C_{01}/132$.
Dokładne przedziały dają $C_{01}>1/50$ i $\|C\|=c_0<1/6$.
Norma śladowa komutatora jest co najmniej podwojonym modułem tego
elementu. W konsekwencji
\begin{equation}\label{eq:explicitdistance}
 \inf_{Q\in\CQ_A}\Dtr(R_{\rm sep},Q)>\frac1{15400}.
\end{equation}
Jest to granica odległości od zbioru klasyczno-kwantowego, nie obliczona
wartość entropowego dysonansu ani miara splątania.

Istnieje też słabsza granica dla wszystkich kanonicznych równowag
z tym $C$. Niech $\widetilde R$ będzie stanem po rzutowaniu na różne
przestrzenie własne $C$ po pierwszej stronie. Z odwracalnego bloku,
stacjonarności i $\|[V,X]\|_1\le2\|V\|\|X\|_1$ otrzymujemy
\[
 \|R-\widetilde R\|_1\ge\frac{c_0-c_6}{2\|V\|\,12}.
\]
Dla najmniejszej niezerowej szczeliny $d_{\min}$ widma $C$,
$\|[R,C\otimes I]\|_F\ge d_{\min}\|R-\widetilde R\|_F$.
Ponieważ pełny wymiar macierzy wynosi 144, norma Frobeniusa jest
co najmniej $1/12$ jej normy śladowej. Łącząc z \eqref{eq:distance},
\begin{equation}\label{eq:universal}
 \inf_{Q\in\CQ_A}\Dtr(R,Q)
 \ge\frac{d_{\min}(c_0-c_6)}{8\|V\|\,12^2(1+c_0)}
 \ge\frac{911271283377980905464649}
          {26400000000000000000000000000000}>0.
\end{equation}
Użyto dokładnych zewnętrznych przedziałów szczelin, a nie
zaokrąglonej diagnozy spektralnej. Dolna granica jest zachowawcza,
około $3.45\cdot10^{-8}$.

Nie należy jej zastępować pozornie mocniejszym argumentem o ustalonym
marginale. Jeśli wymaga się dokładnie tego samego $C$ także od stanu
porównawczego CQ, wykluczenie koincydencji daje dystans co najmniej
$c_0/12$. Ten zbiór porównawczy zmienia się jednak osobliwie przy
$\gamma\to0$: wtedy $C\to I/12$ traci proste widmo, a
$R_{\rm sep}\to(I-D)/132$ jest zwykłym diagonalnym stanem klasycznym.
Nie wolno wnioskować o skończonym dysonansie przy zerowym kodowaniu.

\section{Plan przygotowania nie jest jednoznacznie wyznaczony}

Jedenaście niezorientowanych cięć jest minimalne tylko dla przyjętego
izotropowego drugiego momentu: jego macierz ma rząd 11, a jedno
cięcie wnosi rząd jeden. Nie jest to minimalny koszt wszystkich
możliwych przygotowań kwantowych ani wyprowadzenie fizycznej liczby 12.

Zamiana wierszy 1 i 2 macierzy Hadamarda zachowuje wszystkie statystyki
pierwszego i drugiego rzędu, liczbę cięć, marginały, separowalność
i stacjonarność. Zmienia jednak momenty czwarte. Liczba orientacji,
w których etykiety $0,1$ są po stronie $A$, a $3,4$ po stronie $B$,
rośnie z jednego do dwóch. Dlatego
\[
 (R'_{\rm sep}-R_{\rm sep})_{03,14}
    =\frac2{11}C'_{01}C'_{34}\ne0.
\]
Porównanie całkowitoliczbowych tensorów czwartego momentu we wszystkich
24 przestawieniach dihedralnych wyklucza równoważność tych planów
pod tymi symetriami jądra. Przewidywania dla wspólnej obserwabli
$|03\rangle\langle14|+|14\rangle\langle03|$ są różne.
Jednakowa lokalna geometria nie wybiera więc jednego prawa korelacyjnego.

\section{Próby falsyfikacji i osobny zakres legacy}

\paragraph{Zmiana oddziaływania.}
Twierdzenie o odwracalnym bloku jest kanoniczne, a nie automatycznie
ważne dla wszystkich 4357 parametrów niewidocznych w czystym przepływie.
Przy zastąpieniu $V$ przez $V+aS$ blok \eqref{eq:cross} staje się
osobliwy dla $a=-5/12$. Sam ten test nie daje tam rozstrzygnięcia.
Nie jest to dowód istnienia stanu CQ w przypadku wyjątkowym.
Z kolei czysty SWAP jest inną mapą źródłową i ma stacjonarny produkt
$C\otimes C$ o zerowym dysonansie. Wyklucza to uogólnienie wyniku
na każdą realizację kwantową. Dla perturbacji hamiltonianu o normie
mniejszej niż $1/12$ kanoniczny blok pozostaje odwracalny, o ile
zachowane są te same przesłanki dotyczące marginalnego $C$.

\paragraph{PPT, ranga i rejestry.}
Dowód separowalności jest konstruktywny. Ujemność lub brak ujemności
częściowej transpozycji nie zastępuje tego dowodu.
Poza wystarczającym zakresem wzmocnienia program może przestać być
dodatni; nie oznacza to niemożliwości innej separowalnej konstrukcji.
Również stacjonarność pary nie jest stacjonarnością każdej gałęzi
warunkowanej zachowanym rekordem przygotowania.

\paragraph{Legacy.}
Dla osobnego kanonicznego cyklu
$W_d=4\ln2\cos(\pi d/4+\pi/6)/(1+0.01d)$,
kodowanie $1/1000$ wymaga programu o kodowaniu $11/5000$.
Oszacowanie $|W_{ij}|<3$ daje dodatnią dolną granicę
$1/12-33(11/5000)$, więc ta sama konstrukcja separowalna działa.
Prostotę i nierówność modów 0 i 6 sprawdzono dokładnie w
$\mathbb Q(\sqrt2,\sqrt3)$ po wyłączeniu wspólnego czynnika $4\ln2$.
Wynik jakościowy o niezerowym dysonansie stosuje się zatem osobno
także do tej referencji. Liczb \eqref{eq:explicitdistance}--\eqref{eq:universal}
nie przenosimy na referencję legacy o wagach zmiennego znaku ani na jej role fizyczne.

\section{Bilans i najważniejszy dalszy problem}

\paragraph{Udowodniono.}
Wymagane korelacje dla kanonicznego strict mogą być separowalne;
jawny rozkład i instrument przygotowania spełniają poprzednie
równanie kompensacji komutatora. Jednocześnie każda równowaga
kanonicznego oddziaływania z rzeczywistym marginalnym $C$ z \eqref{eq:C}
ma niezerowy jednostronny dysonans. To precyzyjne rozdzielenie
splątania, separowalności i pełnej klasyczności jest głównym wynikiem.

\paragraph{Obalono lub ograniczono.}
Nie można utożsamiać wymaganych korelacji ze splątaniem, uznawać PPT
za wystarczający certyfikat separowalności ani nazywać wyjścia
przygotowania z klasyczną komunikacją automatycznie stanem CQ.
Minimalna liczba cięć i ich statystyki par nie ustalają pełnej
korelacji. Nie wolno też przenosić kanonicznego no-go na dowolne
zmienione oddziaływanie.

\paragraph{Pozostaje warunkowe.}
Program o większej koherencji, prawo losowania cięć, rekord sukcesu
i operacyjne znaczenie etykiet są nadal dostarczane. Mogą być
legalnymi zasobami fizycznego modelu, ale sama ich konstrukcja nie
wyprowadza ich z FIN. Nie uzyskano źródła jednostek, aparatu,
zapisu laboratoryjnego, selektora QW-2191, mostu legacy--strict,
transferu ról, SM/GR, $L_{\rm total}$ ani ToE. Nadsoliton pozostaje
pierwotną informacją w stanie solitonicznym, bez osobnej warstwy pod nim.

Najbardziej rozstrzygający następny kierunek to niezależne uzasadnienie
stanu programu i jego korelacyjnego prawa wyższych momentów, a następnie
sprawdzenie przewidywań wspólnych obserwabli. Nie wystarczy ponownie
dopasować lokalnego jądra. Trzeba ustalić, czy wymagany zasób koherencji
i wybór przygotowania są rzeczywistymi konsekwencjami teorii, czy jej
empirycznie dostarczanymi warunkami początkowymi.

\section*{Źródła i dane sprawdzalne}
\begin{enumerate}
\item H.~Ollivier, W.~H.~Zurek, \emph{Quantum Discord: A Measure of the
Quantumness of Correlations}, Physical Review Letters 88, 017901 (2001),
\href{https://doi.org/10.1103/PhysRevLett.88.017901}{tekst źródłowy}.
\item C.~A.~Rodríguez-Rosario, G.~Kimura, H.~Imai, A.~Aspuru-Guzik,
\emph{Lazy states: sufficient and necessary condition for zero quantum
entropy rates under any coupling to the environment},
\href{https://arxiv.org/abs/1004.5405}{arXiv:1004.5405}.
\item Wcześniejsze raporty repozytoryjne wskazane w pierwszym przypisie
i dokładne certyfikaty widma. Ich wyniki nie są ponownie liczone jako
odkrycia obecnego opracowania.
\item \texttt{fin\_separable\_stationarity/research.py},
\texttt{discord.py}, \texttt{results.json}, \texttt{discord\_results.json}
i testy. Rozkład produktów oraz rachunek całkowitoliczbowy i wymierny
stanowią certyfikaty; reszty zmiennoprzecinkowe służą oddzielnym kontrolom.
\end{enumerate}
\end{document}
